\documentclass[a4paper,12pt]{article}
\usepackage[finnish]{babel}
\usepackage[utf8]{inputenc}
\usepackage[T1]{fontenc}
\usepackage{lmodern}
\usepackage{eurosym}
\usepackage[pdftex,colorlinks]{hyperref}
\renewcommand{\thesubsection}{\arabic{subsection}}
\title{Etelä-Espoon sosialidemokraattinen työväenyhdistys r.y.}
\author{EED}
\date{19.11.2025}
\begin{document}
\maketitle
\tableofcontents
\section*{Johtokunnan kokous}
\subsection*{Paikka ja aika}
Ravintola Tillinkulma, keskiviikkona 19.11.2025 kello 18:00.
\subsection*{Läsnä}
Markku Sistonen, puheenjohtaja \\
Martti Hellström, varapuheenjohtaja \\
Matti Ikonen \\
Arja Kivimäki \\
Anita Paatelma \\
Pauliina Sistonen \\
Helena Tiitinen \\
Mikko Nummelin, sihteeri
\section*{Pöytäkirja}
\subsection{Kokouksen avaus}
Puheenjohtaja avasi kokouksen kello 18:06.
\subsection{Kokouksen laillisuus ja päätösvaltaisuus}
Kokous todettiin lailliseksi ja päätösvaltaiseksi.
\subsection{Vanhojen pöytäkirjojen tarkastaminen}
Pöytäkirjaan 12.3. kohtaan 9 päätettiin kesän 2025 Falklammen mökin hinnat.
Jäsen: viikko 200\euro, viikonloppu 100\euro, arkipäivä 30\euro. SDP: viikko 300\euro, viikonloppu 120\euro. Muut: viikko 400\euro, viikonloppu 150\euro. \\
Pöytäkirjaan 15.5. kohtaan 5 (Jäsenasiat), yhdistykseen liittyivät Juri Smirnov ja Jani Oksanen. Vaalitulosten käsittelykohdassa käsiteltiin neuvotteluja luottamustehtävien jaosta eri puolueiden välillä. \\
Pöytäkirjaan 3.9. kohtaan tapahtumat, kyseessä oli Demarinaisten 125-vuotisjuhla. \\
Pöytäkirjasta 1.10 puuttui läsnäolijalista. \\
Kaikki edellä mainitut pöytäkirjat hyväksyttiin näin kirjatuin muutoksin.
\subsection{Kokouksen ääntenlaskijoiden valinta}
Kokouksen ääntenlaskijoiksi valittiin Martti Hellström ja Pauliina Sistonen.
\subsection{Taloustilanne}
Todettiin taloustilanne.
\subsection{Jäsenasiat}
Jäsenistössä ei ollut muutoksia.
\subsection{Syksyn tapahtumat}
\begin{itemize}
\item{Joululaulutilaisuus palvelutalossa (ei toteudu)}
\item{to 27.11. kello 18: Puurojuhla (Ravintola Tillinkulma)}
\end{itemize}
\subsection{Muut asiat}
Yhdistys piti jäsenkyselyn jäsenkirjeessä. Toinen kysymyksistä koski matkoja, toinen sähköpostitse saatavaa jäsenkirjettä. 8 jäsentä 75:stä halusivat jäsenkirjeen sähköpostitse. Matkojen osalta tuli 5 vastausta, joissa esitettiin toivomuksia etelänmatkoista lämpimiin maihin, lähialueelle suuntautuvia kaupunkimatkoja sekä kulttuuritapahtumia ja matkoja, joilla on SDP:n aatteellista sisältöä. Yksittäisessä vastauksessa pitkiä ulkomaanmatkoja arvosteltiin tarpeettomina.
\subsection{Kokouksen päättäminen}
Puheenjohtaja päätti kokouksen 19:27.
\end{document}
